% Generated by code/13_generate_exhibits.R
\begin{table}[H]
\centering
\scriptsize
\setlength{\tabcolsep}{2pt}
\caption[Policy-facing country typology]{Policy-facing country typology from PCA-informed clustering}
\label{tab:country_typology_clusters}
\begin{adjustbox}{max width=\textwidth,center}
\begin{tabular}{@{}L{0.85cm}L{3.35cm}rrrrrrr@{}}
\toprule
Type & Countries & $n$ & Rev. & Rev./CIT & QDMTT & Res. & SbS ext. & Imp. \\
\midrule
T1 & LUX &  1 &  57.5 & 211.1\% & 90.1\% &  9.9\% & 28.2\% & 42.1\% \\
T2 & AUT, BEL, DNK, FIN, GRC, HUN, IRL, NLD, SVK &  9 & 173.8 &  15.6\% & 82.4\% & 17.6\% & 23.9\% & 25.0\% \\
T3 & BGR, CYP, CZE, DEU, ESP, FRA, HRV, ITA, POL, PRT, ROU, SVN, SWE & 13 & 136.3 &   6.3\% & 44.7\% & 55.3\% & 39.0\% & 34.9\% \\
T4 & EST, LTU, LVA, MLT &  4 &   0.0 &   0.0\% &  &  &  & 93.5\% \\
\bottomrule
\end{tabular}
\end{adjustbox}

\vspace{0.35em}
\noindent\begin{minipage}{\textwidth}
\footnotesize\raggedright\textit{Notes:} Four clusters are selected because the jitter stability score is 100.0\%, the smallest cluster contains one country, and the four-cluster solution separates conduit hubs (T1), QDMTT-protected smaller states (T2), residual IIR/UTPR exposure states (T3) and states with zero baseline collector receipts (T4). Rev.\ is preferred behavioural SbS revenue in US dollars, billions. T4 has zero assigned receipts and undefined channel shares. Estonia, Latvia, Lithuania and Malta deferred IIR and UTPR under Article 50. Rev./CIT uses OECD T\_1200 where available; cluster means ignore missing denominators. Res.\ is residual IIR/UTPR exposure; SbS ext.\ is the share of CIT revenue extinguished under the SbS regime; Imp.\ is the high-risk imputation share.
\end{minipage}
\end{table}
